\documentclass{article}
\usepackage{iclr2026_conference,times}
\input{math_commands.tex}
\usepackage{graphicx}
\usepackage{booktabs}
\usepackage{hyperref}
\usepackage{url}

\title{Beyond Preliminary INT4 Results: A Deeper Study of Module-Aware Quantization for DINO-WM Planning}
\author{}

\begin{document}
\maketitle
\lhead{Under review as a workshop paper at ICLR 2026 (ES-Reasoning)}

\begin{abstract}
We revisit the preliminary finding that mixed INT4 quantization (encoder FP16, predictor INT4) outperforms uniform INT4 for DINO-WM planning on Wall.
Our goal is to test whether this effect is robust, statistically stable, and mechanistically tied to where precision is allocated.
We expand the original setup along six research questions: multi-budget robustness of mixed-vs-uniform INT4, closeness of mixed INT4 to INT8/FP16, position-sensitive encoder retention, asymmetric encoder/predictor bit allocation, random-goal robustness, and episode-level paired matchup dominance.
This draft is synchronized with the M3 experiment pipeline and is intended for final number insertion after run completion.
\end{abstract}

\section{Introduction}
Low-bit quantization in world models is often framed as a one-dimensional question of bitwidth.
Our prior result suggests a stronger hypothesis: near the 4-bit transition regime, \emph{where} bits are preserved can matter more than average bitwidth alone.
This paper deepens that hypothesis beyond preliminary evidence by running targeted robustness and allocation-structure experiments.

\section{Research Questions}
\begin{itemize}
  \item \textbf{RQ1}: Is mixed INT4 consistently better than uniform INT4 across planning budgets and seeds?
  \item \textbf{RQ2}: How close is mixed INT4 to INT8 and FP16 performance under matched planning budgets?
  \item \textbf{RQ3}: Inside the encoder, does preserving tail blocks vs head blocks in FP16 change INT4 outcomes?
  \item \textbf{RQ4}: At similar average precision, does allocation direction (encoder-heavy vs predictor-heavy) alter success?
  \item \textbf{RQ5}: Does the mixed INT4 gain persist when goal source changes from paired files to random-state goals?
  \item \textbf{RQ6}: At the episode level, does mixed INT4 win on the same paired episodes or only in aggregate means?
\end{itemize}

\section{Experimental Protocol}
We evaluate DINO-WM on Wall using paired-goal and random-goal budgets, with M3 Apple Silicon (MPS backend).
The expanded variant families include:
\begin{itemize}
  \item baseline/frontier: FP16, uniform/mixed INT8, uniform/mixed INT6, uniform/mixed INT4;
  \item encoder position studies: tail-preserved FP16 (encfp16\_25/50/75) vs head-preserved FP16 (encheadfp16\_25/50);
  \item asymmetric mixed families: E8/P4 vs E4/P8 and E6/P4 vs E4/P6.
\end{itemize}
Primary statistics are paired episode deltas with bootstrap confidence intervals, augmented by episode-level contingency counts.

\section{Results}
\subsection{RQ1: Mixed INT4 vs Uniform INT4}
Insert paired deltas for budgets bA/bB/bC:
\begin{itemize}
  \item \texttt{paired\_delta\_bA\_mixed\_vs\_uniform\_int4.json}
  \item \texttt{paired\_delta\_bB\_mixed\_vs\_uniform\_int4.json}
  \item \texttt{paired\_delta\_bC\_mixed\_vs\_uniform\_int4.json}
\end{itemize}

\subsection{RQ2: Closeness to INT8/FP16}
Report absolute success gaps of mixed INT4 vs mixed INT8, uniform INT8, and FP16 from grouped summaries.

\subsection{RQ3-RQ4: Allocation Structure}
Report tail-vs-head paired deltas and cross-allocation deltas (E8/P4 vs E4/P8, E6/P4 vs E4/P6).

\subsection{RQ5: Goal Source Robustness}
Insert random-goal deltas:
\begin{itemize}
  \item \texttt{paired\_delta\_bA\_rand\_mixed\_vs\_uniform\_int4.json}
  \item \texttt{paired\_delta\_bB\_rand\_mixed\_vs\_uniform\_int4.json}
\end{itemize}

\subsection{RQ6: Episode-Level Matchups}
Insert contingency results:
\begin{itemize}
  \item \texttt{episode\_matchup\_int4\_mixed\_vs\_uniform.json}
  \item \texttt{episode\_matchup\_tail50\_vs\_head50.json}
\end{itemize}

\section{Discussion}
Focus on whether evidence supports the central claim:
module-aware precision allocation at 4-bit yields stronger planning reliability than uniform 4-bit quantization, and whether this advantage narrows or remains close to INT8/FP16 in key budgets.

\section{Limitations}
Single environment (Wall), one pretrained checkpoint, Apple M3 hardware, and weight-only quantization path.

\section{Conclusion}
State only claims directly supported by paired deltas and episode-level dominance outcomes.

\bibliography{references}
\bibliographystyle{iclr2026_conference}

\end{document}
